\section{Introduction}
\label{sec:intro}

In every model and dataset cell we tested, a language model's intermediate layers scored its own hallucinations better than its output layer did --- using nothing but the uncertainty statistics of the token distributions the model already computes on the way to its answer. That is our first finding. The second concerns measurement rather than models: the documented failure baseline this study was designed to beat did not survive a provenance audit. Six protocol differences --- sample selection, prompt template, label rule, signal definition, numeric precision, and evaluation split --- moved the same cell's final-layer entropy AUROC from 0.5186 to 0.5928 before a single new claim had been tested. We treat the second finding as a result, not an embarrassment.

The surface problem is familiar: open-weight LLMs hallucinate, and the cheapest detection signals --- confidence statistics read from the output distribution --- are architecture-fragile. In an archived record from an earlier study in this research line --- hereafter the \emph{motivating record} --- final-layer mean token entropy reached AUROC 0.66 on LLaMA-3-8B-Instruct but roughly 0.52, with inverted score direction, on base LLaMA-2-7B: the same statistic, the same datasets. We cite these numbers only as motivating context; the audit above shows they are bound to that record's protocol, and they are never used as same-protocol baselines here. The fragility itself, however, is the point. A practitioner cannot ship a confidence score whose validity flips with the checkpoint, and the standard escapes are expensive: multi-sample semantic consistency pays roughly ten forward passes per query \citep{Farquhar2024Detecting}, and supervised hidden-state probes require labeled data and per-model training \citep{Azaria2023Internal}.

The deeper problem is that the field has converged on a diagnosis without adopting the cheap cure. That intermediate layers carry stronger truthfulness signal than the final layer is by now close to consensus \citep{Orgad2024LLMs,Wang2026Automatic}. Yet every method that exploits this re-introduces machinery the observation should have made unnecessary: trained probes over hidden states \citep{Azaria2023Internal,Suresh2025Cross}, geometric layer-selection criteria feeding supervised classifiers \citep{Wang2026Automatic}, or repeated sampling \citep{Chen2024INSIDE}.

This leaves a specific gap. The raw logit lens --- projecting each layer's hidden state through the model's own unembedding --- yields a full next-token distribution at every depth, and the uncertainty statistics of those distributions have never been evaluated as hallucination-detection scores. Adjacent-layer KL divergence, a natural measure of how much the model revises its belief between layers, has been used only to steer decoding \citep{Chuang2023DoLa,Wu2025Improve}, never scored for detection. The gap persisted partly through a community split --- interpretability owns the instrument and studies these statistics as computation signatures \citep{Ali2025Entropy}, while detection owns the task and defaults to probes and sampling --- and partly through discouragement by adjacency: the nearest published measurement tracked final-prediction-token probability trajectories across layers and found certain and uncertain trajectories largely aligned \citep{Kim2025Effect}, a negative result for a different statistic that was easy to over-read as covering this one.

Our premise is that the gap hides signal because of where the standard signal is read. The final layer is not where the information dies; it is where output calibration --- a tokenizer- and tuning-dependent presentation step --- reshapes it. Watch the model make up its mind layer by layer: when it knows the answer, next-token candidate competition collapses at intermediate depth and stays collapsed; when it is hallucinating, the decoded distribution stays wide and keeps churning until the output head tidies it up. The design principle that follows is to read before calibration: compute entropy, max-token probability, and adjacent-layer KL from every layer's logit-lens distribution --- obtained from one greedy generation plus one teacher-forced re-forward, with no sampling and no training --- and select the readout layer, signal, and score direction per model on a held-out selection split, because calibration depth is a property of the checkpoint.

We test this premise at proof-of-concept scale under a staged verification design whose first tier asks only whether the signal exists. Across LLaMA-2-7B, Mistral-7B-v0.1, and LLaMA-3-8B-Instruct on TriviaQA and TruthfulQA --- 5,451 scored generations covering all six model $\times$ dataset cells --- at least one screened intermediate layer is class-separable in every cell, with selection-split corrected AUROC between 0.6092 and 0.7011 against a 0.55 existence gate, and the best screened intermediate layer beats the same sweep's final-layer entropy readout in every cell (point-estimate margins +0.0035 to +0.130). These are selection-split point estimates from a single seed, with no confidence intervals: the numbers that chose the (layer, signal, direction) tuples also grade them, and they must not be read as held-out performance. The locked test splits exist and were never touched --- by any analysis, or by any decision made in writing this paper --- which preserves the deferred evaluation as an uncontaminated, pre-registered confirmation; running the frozen tuples there, along with cross-dataset transfer and KL--entropy fusion, is the explicitly unmeasured next tier of the staged design. Within this scope, the depth advantage is universal in direction and consistent with the calibration-suppression reading, though it does not establish that mechanism.

Four contributions follow from the insight. First, the first AUROC evaluation of raw per-layer logit-lens uncertainty statistics as training-free, sampling-free hallucination scores: class-separability holds in all six cells at existence tier. Second, a universal point-estimate depth advantage: the best screened intermediate layer --- concentrated at L28--L31 --- exceeds the within-sweep final-layer baseline in every cell, extending the intermediate-over-final consensus to statistics that need no probe. Third, adjacent-layer KL as a detection signal in its own right: previously a decoding heuristic, it is the best signal in two of our six cells --- an even three-way split with entropy and max-probability --- and sweeps the top three intermediate scores in the binding LLaMA-2/TriviaQA cell. Fourth, a methodological contribution born from failure: protocol-internal validity anchoring. Having shown quantitatively that AUROC magnitudes for weak signals do not survive protocol changes --- direction transfers, magnitude does not --- we replace cross-run numeric anchors with identity-verified cache reuse, within-sweep baselines, and descriptive-only cross-run comparisons.

This work sits at a junction that, to our knowledge, has never been connected without supervision: the interpretability community's per-layer lens instruments on one side, and hallucination detection on the other. We next situate it among the lines of work that meet there.
